\section{Eclector: A highly configurable reader}
\label{sec-eclector}

\subsection{Features}

We could talk about the following unusual features:

\begin{itemize}
\item Production of parse results with source information.  Including
  skipped input (comments, reader conditionals) and invisible features
  of the input (labeled object definitions and references, feature
  expressions in reader conditional)

\item Specific conditions for basically every possible syntax error.

\item Recovery from every possible syntax error.

\item Most aspect of the reader are customizable. Customization in
  standard \cl{}: Reader control variables and Readtable (case, syntax
  type, macros)

\item Intrinsic and extrinsic usage

\item Reader state protocol

\end{itemize}

\subsection{Robustness and Unusual Inputs}

We can talk about the ``hybrid recursion'' strategy for handling
deeply nested structures and similar issues.  We could give examples
that, say, \sbcl{} cannot handle.

\subsection{Efficiency}

We could do a few benchmarks

Another thing we could talk about are edge cases of labeled objects
and references.  Not because the performance of those is particularly
relevant in practice but because there is algorithmic innovation in
\eclector{} in this area and we even have an edge over \sbcl{} for
certain cases.
